\section{From actual level modules to the full Taylor--Wiles automorphy datum}
\label{sec:hida-taylor-wiles-v45}

This chapter proves the finite consumers between the original
\(\mathrm{LevelAlt}\) and \(\mathrm{HasAut}\).  Its only new
representation-theoretic premises are precisely the three fields
of the original \(\mathrm{LocPAlt}\): Bernstein--Zelevinsky
finiteness, the geometric lemma, and the tame local theorem.
The full \(\mathrm{IsRecT}\) correspondence laws, its existence,
and its generic and spherical good-parameter consequences have
the scope already reconstructed in GUA.  The actual level module
is supplied by HLM, with its vector-separating classical maps,
full polynomial-law comparisons, and one fixed base module.

Here \(\Lambda=O[[Z_1,\ldots,Z_D]]\), where the empty set of
additional variables is understood and \(D=n[L^+:\mathbb Q]\).
It is a complete noetherian local ring.  All modules described as
finite below are finite over this ring unless another ring is
specified.  The fixed base is \(b\), its module is \(M_0=b.M\),
and its coefficient map is \(\lambda:\Lambda\to B\).  In
particular, for every prime \(\mathfrak p\) of \(B\),
\[
 \operatorname{Ann}_B(M_0)\subseteq\mathfrak p
       \quad\Longleftrightarrow\quad J\subseteq\mathfrak p .
 \tag{HTW.1}
\]
This equivalence is retained throughout; none of the intermediate
localizations changes the chosen ideal \(J\).

At \(q\) auxiliary places write \(Q_j>1\) for the residue
cardinality, \(f_j\) for geometric Frobenius, and \(t_j\) for
the original tame generator.  Indices \(i\) run from \(0\) to
\(n-1\).  For commuting variables \(U_{ji}\) put
\[
 F(U)=\prod_j\prod_i\prod_{k\ne i}
             (U_{ji}-U_{jk})(U_{ji}-Q_jU_{jk}).
 \tag{HTW.2}
\]
This is the literal source polynomial \(\mathrm{fEl}\), with
ordered pairs, not its square-free part.  In the group algebra
define
\[
 \begin{split}
 C_j&=\prod_i([f_j]-U_{ji}),\\
 C_{ji}'&=\left(\prod_{k\ne i}([f_j]-U_{jk})\right)
                                      ([t_j]-1-Y_{ji}).
 \end{split}
 \tag{HTW.3}
\]
The order displayed in the second expression is retained.

\subsection{Finite tame level and the two local relations}

% BEGIN HTW1
\paragraph{HTW1. The tame quotient and its exact norm sequence.}
\label{claim:HTW1}
Let \(I_j\) be the original Iwahori subgroup and \(I_j^1\)
its subgroup with diagonal reductions in the prime-to-\(\ell\)
subgroup of the residue-field units.  If
\(a_j=v_\ell(Q_j-1)\), then
\[
 \Delta=\prod_j I_j/I_j^1
       \simeq\prod_{j,i}\mathbb Z/\ell^{a_j}\mathbb Z .
 \tag{HTW.4}
\]
The original \(\delta_j\) gives the generator in each diagonal
coordinate.  For any projective \(\Lambda[\Delta]\)-module \(M\),
the norm \(N_\Delta=\sum_{\gamma\in\Delta}\gamma\) induces
an isomorphism \(M_\Delta\simeq M^\Delta\).  Its kernel on \(M\)
is \(\sum_{j,i}(\gamma_{ji}-1)M\).
\begin{proof}
Reduction of the diagonal entries is surjective onto the
product of residue-field unit groups; its kernel contains the
unipotent factors and the principal congruence group.  Taking
the specified prime-to-\(\ell\) diagonal subgroup leaves the
cyclic \(\ell\)-primary quotient in each coordinate.  A finite
field's multiplicative group is cyclic, and a generator of
its quotient has order \(\ell^{a_j}\).  The defining
\(\mathrm{IsGenMod}\) property of \(\delta_j\) says exactly
that its class is such a generator.  Different places and
different diagonal coordinates commute in this quotient.

For the regular module \(\Lambda[\Delta]\), a fixed vector
has all coefficients equal and is uniquely a scalar multiple
of \(N_\Delta\).  The norm of \(\sum c_\gamma[\gamma]\)
is \((\sum c_\gamma)N_\Delta\).  Its kernel is the
augmentation ideal: subtract the coefficient sum at \(1\)
to express it as \(\sum c_\gamma([\gamma]-1)\).
Writing every \(\gamma\) as a product of the indicated
generators and using \(ab-1=a(b-1)+(a-1)\) proves the
generator description.  These assertions hold for free
modules coordinatewise and pass to a direct summand.
This proves the projective case without division by
\(|\Delta|\).  Finally \(Q_j\equiv1\bmod\ell^N\) implies
\(a_j\ge N\), also when \(N=0\).
\end{proof}
% END HTW1

% BEGIN HTW2
\paragraph{HTW2. The actual commuting torus operators and trace.}
\label{claim:HTW2}
For an original level module \(A\), put
\(M_1=A.P^{\prod I_j^1}\) and \(W=A.P^{\prod I_j}\).
There are commuting invertible torus actions \(\Theta_1,\Theta_w\)
with the original dominant double-coset values.  Define
\[
 U_{ji}=Q_j^i\Theta_1(e_i(x_j)),\qquad
 1+Y_{ji}=\Theta_1(e_i(\delta_j)),
 \tag{HTW.5}
\]
and use the same \(U\)-formula on \(W\).  The actual
double-coset trace \(M_1\to W\) is the norm in HTW1.
It is surjective, its kernel is \(\sum Y_{ji}M_1\), and
it intertwines every \(U_{ji}\).
\begin{proof}
All \(Q_j\) are units in \(O\), since the places avoid
\(\ell\).  The positive diagonal double cosets multiply
according to addition of their ordered valuation vectors.
This follows by writing an Iwahori element in upper
unipotent, diagonal-unit, and lower unipotent order:
the summation indices for a product are the disjoint
successive residue quotients of each root subgroup.
At the tame subgroup the simple-reflection relation
retains its diagonal-unit terms.  For a nonzero residue
class \(c\), choose a unit lift \(a_c\) and let \(d_c\)
have entries \(a_c,-a_c^{-1}\) in the two adjacent
coordinates and \(1\) elsewhere.  Write \(D_c\) for
its normalizer action.  The elementary two-by-two
matrix factorization of \(s(1+a_cE_{i,i+1})s\) gives
\[
 T_s^2=Q_j+T_s\sum_{c\ne0}D_c,\qquad
       T_sD_c=D_c^{\,s}T_s .
\]
Consequently \(Q_j^{-1}(T_s-\sum D_c)\) is a right
inverse and \(Q_j^{-1}(T_s-\sum D_c^{\,s})\) is a
left inverse.  They agree, so \(T_s\) is a unit.
Only on the ordinary Iwahori fixed space do all
\(D_c\) act as \(1\), giving
\(T_s^2=(Q_j-1)T_s+Q_j\).

Reduced permutation products are therefore units at
either level.  For a fundamental diagonal translation
with uniformizer in the first \(m\) entries, rotating
those entries to the end makes its product with the
rotation permutation normalize \(I_C\).  Indeed
conjugation changes the valuation bound on matrix
entry \((i,k)\) by the difference of these two
zero-or-one exponents, exactly compensating the
change of whether the entry is below the diagonal;
it permutes the prescribed diagonal residue subgroup.
The translation double coset is thus a product of
a permutation unit and a normalizer unit.
An ordered integral valuation vector is a sum of
these fundamental vectors and a constant vector.
Positive multiplication, and the normalizer actions
of diagonal units and central scalars, prove that
every dominant diagonal operator is invertible.
Every diagonal element is a quotient of two
dominant ones: add a sufficiently strictly ordered
valuation vector to make it and its translate dominant.
The positive multiplication rule therefore extends
uniquely to the torus.  Applying this independently at
each place gives commuting actions and (HTW.5).

The right cosets in \(I_j/I_j^1\) have diagonal-unit
representatives.  Thus the double coset at the identity
is exactly their sum, with no volume factor.  For a
dominant diagonal \(t\), the two expressions obtained
by first summing these units and then its root-subgroup
representatives, or in the other order, enumerate the
same cosets.  The trace intertwines \(\Theta(t)\).
For an arbitrary \(t\), multiply by a dominant \(b\)
such that \(bt\) is dominant and cancel the invertible
\(\Theta(b)\).  Hence it intertwines all torus operators,
including (HTW.5).  The level-module projectivity field
and HTW1 give the asserted kernel and surjectivity.
\end{proof}
% END HTW2

% BEGIN HTW3
\paragraph{HTW3. The uniform local nilpotence exponent is \(n!\).}
\label{claim:HTW3}
Let \(\sigma\) be smooth irreducible over \(\mathbb C\).
For the original subgroup \(I_C\), its fixed space has
dimension at most \(n!\).  A polynomial in the commuting
torus operators which vanishes at every joint generalized
character acts nilpotently with exponent \(n!\).
The same exponent works in an arbitrary isotypic sum of
copies of \(\sigma\).
\begin{proof}
Bernstein--Zelevinsky finiteness makes the fixed space finite
dimensional.  Its map to the Borel Jacquet quotient is
injective: if a fixed vector is a finite sum of unipotent
differences, a sufficiently dominant diagonal contracts
the finitely many unipotents into the fixing subgroup.
The translated vector is then a sum of differences
for elements of that subgroup.  Every root-coset
translate remains in this difference span, so the
positive double-coset operator sends the vector to
both that span and the fixed space.  Normalized
compact averaging kills the former and fixes the
latter; it is a finite average on the finitely
many smooth vectors involved.  Thus this operator
kills the vector.  HTW2 makes it invertible, so
the vector was zero.

If the fixed space is nonzero, commuting invertible torus
operators have a common eigenvector.  To see this, choose
a minimal nonzero invariant subspace; the eigenspace of
any one operator inside it is invariant under all the
others, so minimality makes every operator scalar there.
The image in the Jacquet module is nonzero.
The retained finite-length theorem for that module,
and the fact that smooth irreducible torus modules
are characters, give a finite-dimensional torus
module.  The latter fact follows by an open
stabilizer and the weak Nullstellensatz for the
remaining finitely generated abelian group, as
in HAW7.  Choose a nonzero generalized-character
block.  Its quotient by the maximal ideal of the
commutative image algebra is nonzero by Nakayama;
a nonzero functional on that quotient gives a
character quotient of the entire Jacquet module.
It need not detect the initial eigenvector.
Jacquet reciprocity applied to this quotient
gives a nonzero map from \(\sigma\) into the
corresponding principal series, which is injective
by irreducibility.  The geometric lemma gives the Jacquet
module of that principal series a filtration with one
one-dimensional character for each permutation.  Exactness
of Jacquet coinvariants gives the required bound \(n!\).
This uses the original geometric lemma on this principal
series, rather than any bound on an isotypic multiplicity.

A commuting family on a finite-dimensional complex space
can be simultaneously triangularized, by the same common
eigenvector argument and induction on the quotient.  Its
diagonal characters occur among the joint generalized
characters: equivalently decompose the finite-dimensional
commutative image algebra into its local Artinian factors.
The polynomial in the assertion therefore has zero
diagonal.  A strictly upper triangular matrix on a space
of dimension \(d\le n!\) has \(n!\)-th power zero.

Finally every fixed vector in an isotypic sum is a finite
sum of fixed vectors from actual copies of \(\sigma\).
Indeed start with a finite expression in the images of
intertwiners and average over the compact fixing subgroup;
smoothness makes each average a finite sum.  Intertwiners
commute with the torus action on dominant elements by the
finite double-coset formula, and on all elements by
HTW2's cancellation.  The same polynomial power kills
every summand.  Its exponent does not increase with the
number of copies or the number of auxiliary places.
\end{proof}
% END HTW3

% BEGIN HTW4
\paragraph{HTW4. Tame-kernel triviality from the actual Weil--Deligne comparison.}
\label{claim:HTW4}
Away from \(\ell\), an actual \(\mathrm{IsWDOf}\) comparison
gives equality of characteristic polynomials on the Weil
group.  An element of the kernel of the pro-\(\ell\) tame
character whose matrix has characteristic polynomial
\((X-1)^n\) acts as the identity.
\begin{proof}
On a sufficiently small open inertia subgroup, the
monodromy formula is
\[
                  \rho(\tau)=\exp(t_\ell(\tau)N).
\]
On the
specified tame kernel the parameter \(t_\ell\) is zero.
Some positive power of the given inertia element lies in
that open subgroup, because inertia is compact and the
quotient is finite.  Its matrix therefore has finite
order \(d\).  Its minimal polynomial divides both
\(X^d-1\) and \((X-1)^n\).  The former is separable in
characteristic zero, so the minimal polynomial divides
\(X-1\), and the matrix is \(1\).

For the first assertion, filter by the kernels of
successive powers of \(N\).  The raw Weil operator
preserves this filtration because it conjugates \(N\)
to a nonzero scalar multiple of \(N\); its exponential
correction is the identity on each successive quotient.
The corrected and uncorrected operators therefore
have identical characteristic polynomials.  The
Frobenius semisimplification also preserves these
polynomials, as in the full \(\mathrm{IsWDOf}\)
reading reconstructed in GUA.  Thus the result is the
Weil--Deligne characteristic polynomial at every Weil
element, with the same geometric Frobenius convention.
Coefficient transport uses only the field isomorphism
\(\iota\), never its continuity.
\end{proof}
% END HTW4

% BEGIN HTW5
\paragraph{HTW5. The nonresonant tame matrix identity.}
\label{claim:HTW5}
Suppose a representation \(R\) of a group has, for every \(g\),
\[
 \det(X-R(g))=\prod_i(X-c_i(g)),
\]
where the \(c_i\) are characters.  Put \(P=R(f)\), \(T=R(t)\),
\(\alpha_i=c_i(f)\), \(\beta_i=c_i(t)\).  Assume
\(P^{-1}TP=T^Q\), \(Q\ge2\), and
\(\alpha_i\ne\alpha_k,\ Q\alpha_k\) for \(i\ne k\).
Then
\[
 \prod_{k\ne i}(P-\alpha_k)\,(T-\beta_i)=0 .
 \tag{HTW.6}
\]
\begin{proof}
The character relation gives \(\beta_i^{Q-1}=1\).
Thus \(T^{Q-1}\) is unipotent by its characteristic
polynomial.  Let \(L=\log(T^{Q-1})\), using the finite
nilpotent logarithm over \(\mathbb C\).  The tame relation
gives \(P^{-1}LP=QL\).  The distinct eigenvalues of \(P\)
give an eigenbasis; in it each entry of \(L\) satisfies
\((\alpha_k/\alpha_i-Q)L_{ik}=0\).  All these coefficients
are nonzero: for different indices this is the stated
ordered-pair condition, and for equal indices it is
\(1-Q\).  Consequently \(L=0\), hence \(T^{Q-1}=1\).
The tame relation now gives \(PT=TP\).

In that eigenbasis \(T\) is diagonal, say with entries
\(d_i\).  For \(a=0,\ldots,n-1\), the trace identity at
\(f^at\) says
\(\sum_i\alpha_i^a d_i=\sum_i\alpha_i^a\beta_i\).
The Vandermonde determinant is nonzero, so \(d_i=\beta_i\).
The product in (HTW.6) is zero on every eigenline: its
first factor kills all lines except the \(i\)-th, and its
last factor kills that line.  This proves the ordered
matrix identity.  Rank zero is the unique zero-dimensional
matrix identity and needs no eigenbasis.
\end{proof}
% END HTW5

% BEGIN HTW6
\paragraph{HTW6. The full original local-data consumer.}
\label{claim:HTW6}
The original \(\mathrm{LocPAlt}\) and the full local comparison
at every classical point produce the literal
\(\mathrm{LocalLevelConditions}\), with the given geometric
Frobenius and tame generators.
\begin{proof}
Fix a prolongation and one permissible \(\mathrm{IsRecT}\)
correspondence, whose existence is among the retained
original laws.  The tame local theorem and the original
Artin map send \(f\) to a uniformizer \(x\), send \(t\)
to a unit \(\delta\) generating the indicated residue
quotient, and put \(t^{-Q}f^{-1}tf\) in the pro-\(\ell\)
tame kernel.  This choice is made before the rank,
classical point, irreducible component, or its torus
character is chosen.  After these choices fix an
arbitrary residue-unit subgroup \(C\supseteq C_\ell\),
an arbitrary admissible \(\Theta\) on the corresponding
\(I_C\)-fixed space, and any occurring joint
generalized character \(\chi\).

A joint generalized torus character \(\chi\) on the
component gives, by Jacquet reciprocity, smooth inducing
characters \(\psi_i\), trivial on the prescribed unit
subgroup, with
\[
 \psi_i(z)=Q^{i\,\operatorname{ord}(z)}
     \chi(e_i(z))\,|z|^{(n-1)/2}.
 \tag{HTW.7}
\]
This follows by evaluating the half-modulus character
of the upper Borel at the single diagonal coordinate:
its exponent is \((n-1)/2-i\).
Set \(\alpha_i=Q^i\chi(e_i(x))\).
If all the factors in (HTW.2) at this place are nonzero,
no two inducing characters are related by the modulus
character: evaluating any such relation at \(x\) would
give \(\alpha_k=Q\alpha_i\).  Thus the toral support is
unlinked.  GUA2's full \(\mathrm{PlainGd}\) law applies
to \(\mathrm{algNorm}(\sigma)\).  A spherical component
also has a good parameter by GUA12.
These are exactly the two permitted alternatives in
the local-data contract.  Only after proving this
good-parameter guard do we use the point's conditional
\(\mathrm{WDT}\) comparison to obtain the actual
\(\mathrm{IsWDOf}\) witness.

On the Weil group the corresponding characters after
the algebraic normalization are
\[
 c_i(g)=Q^{i\,\operatorname{ord}(a(g))}
                         \chi(e_i(a(g))).
 \tag{HTW.8}
\]
The full local comparison and HTW4 give the characteristic
polynomial \(\prod_i(X-c_i(g))\) for the actual matrix.
At \(f\) these values are \(\alpha_i\); at \(t\) they
are \(\chi(e_i(\delta))\).  At \(t^{-Q}f^{-1}tf\)
all the characters equal \(1\), by their triviality
on the specified unit subgroup.  HTW4 kills the actual
matrix of this element.  HTW5 gives (HTW.6).
In the spherical alternative only the characteristic
polynomial assertion is needed; in the nonresonant
alternative both assertions hold.

The argument uses the classical point's local comparison
for every permissible correspondence and every
prolongation.  Therefore it applies to the selected one
without an extra supplier.  Taking the already supplied
vector-separating classical family completes precisely
\(\mathrm{LocalLevelConditions}\).
\end{proof}
% END HTW6

% BEGIN HTW7
\paragraph{HTW7. Descent of polynomial nilpotence through classical separation.}
\label{claim:HTW7}
Let \(H\in R[U_i,Y_i]\) concern one auxiliary place.
Suppose every classical specialization and every occurring
joint generalized character makes it zero after the
substitution
\[
 U_i=Q^i\chi(e_i(x)),\qquad
 Y_i=\chi(e_i(\delta))-1 .
 \tag{HTW.9}
\]
Then \(H^{n!}\) acts as zero on \(M_1\).
\begin{proof}
Each classical map is additive, coefficient-semilinear,
and equivariant for the actual group action.  Finite
double-coset summation therefore intertwines every
dominant torus operator.  HTW2 extends this to all torus
elements, and polynomial induction extends it to \(H\).
Apply HTW3 to the classical target, which is an isotypic
sum at the chosen place.  Formula (HTW.9) is exactly
the eigenvalue substitution in (HTW.5); hence the
specialization of \(H^{n!}m\) is zero for every point.
The \(\mathrm{ClassicalAlt}\) field separates actual
module vectors, so \(H^{n!}m=0\).  The conclusion is
integral over the unchanged \(R\).
\end{proof}
% END HTW7

% BEGIN HTW8
\paragraph{HTW8. Both exact trace relations with the same exponent.}
\label{claim:HTW8}
For every \(j,i,g\), on \(M_1\) one has
\[
 \big(F(U)D.\operatorname{tr}([g]C_j)\big)^{n!}=0,\qquad
 \big(F(U)D.\operatorname{tr}([g]C'_{ji})\big)^{n!}=0 .
 \tag{HTW.10}
\]
\begin{proof}
First use only the factor \(F_j\) at place \(j\).
At a joint character either \(F_j=0\), making the
required scalar zero, or HTW6 gives the characteristic
polynomial of \(P\) and the matrix identity (HTW.6).
Cayley--Hamilton then kills \(C_j(P)\); (HTW.6) kills
the ordered \(C'_{ji}(P,T)\).  Left multiplication
by the arbitrary matrix \(R(g)\) and taking the trace
preserves zero.  The point's equality of full determinant
laws specializes their trace polynomial to exactly this
matrix trace, over the polynomial scalar algebra.
HTW7 proves both \(n!\)-th power relations for \(F_j\).

The remaining factors of \(F\) commute with these
operators, so multiplying the zero power by their
\(n!\)-th powers gives (HTW.10).  There is one
exponent \(n!\), rather than an exponent multiplied by
the number of places.  In rank zero the degree-zero
determinant has zero trace and both assertions hold
directly.  No division by \(n!\) occurs.
\end{proof}
% END HTW8

% BEGIN HTW9
\paragraph{HTW9. The symmetric operator identities on the original base module.}
\label{claim:HTW9}
Under the inclusion \(M_0\to W\), the elementary symmetric
polynomials in \(U_{j0},\ldots,U_{j,n-1}\) act by
\((-1)^kP_j{}_{\,n-k}\), where
\[
 P_j(X)=D_0.\chi([h(f_j)])\in B[X].
 \tag{HTW.11}
\]
\begin{proof}
The original level-module equivalence identifies \(M_0\)
with the vectors fixed by all maximal integral compact
subgroups.  A nonzero image of such a vector at a
classical point lies in a spherical component.
HTW6 gives its Frobenius polynomial with roots
\(Q_j^i\chi(e_i(x_j))\).  The elementary symmetric
functions in those roots are the coefficients in
(HTW.11) with sign \((-1)^k\).

The torus action on the spherical vector obeys the same
identity by the normalized Satake calculation already
given in BCS15 and UPA5--UPA7; equivalently it follows
by evaluating the elementary symmetric polynomial in
the Iwahori character decomposition.  The Hecke and
determinant scalar actions agree at the same classical
point by its full determinant-law field.  Vector
separation descends the equality to the actual fixed
module.  This holds for \(0\le k\le n\), with monic
degree \(n\) polynomials, and for the original
coefficient map \(\beta:R\to B\).
\end{proof}
% END HTW9

\subsection{Integral comparison maps and the exact correction factor}

% BEGIN HTW10
\paragraph{HTW10. An explicit integral divided-difference dual family.}
\label{claim:HTW10}
In \(O[X_0,\ldots,X_{n-1}]\) set
\[
 d=\prod_{i<k}(X_k-X_i),\quad
 N_Q=\prod_{i<k}(X_k-QX_i),\quad
 \partial f=d^{-1}\sum_{w\in S_n}\operatorname{sgn}(w)w(f).
 \tag{HTW.12}
\]
The last expression is a polynomial.  There are explicit
finite families \(a_T,b_T\), defined over \(\mathbb Z\),
such that
\[
 \sum_T a_Tw(b_T)=
     \begin{cases}d&w=1,\\0&w\ne1,\end{cases}
 \qquad
 \sum_T a_T\partial(b_Tf)=f .
 \tag{HTW.13}
\]
The same expansion holds with \(a,b\) interchanged.
\begin{proof}
An alternating polynomial vanishes on \(X_i=X_k\);
successive division by the monic differences proves
divisibility by \(d\), first over the integral universal
polynomial ring and then after coefficient substitution.
Let \(\mathcal P=\{(i,k):i<k\}\), and index by its subsets.
Put
\[
 a_T=\prod_{(i,k)\in T}X_k,\qquad
 b_T=\prod_{(i,k)\in\mathcal P\setminus T}(-X_i).
\]
Expanding a product gives
\(\sum_Ta_Tw(b_T)=\prod_{i<k}(X_k-X_{w(i)})\).
For \(w=1\) this is \(d\).
For \(w\ne1\), some \(i<w(i)\); otherwise summing
\(w(i)\le i\) forces equality everywhere.  The pair
\((i,w(i))\) makes the product zero.
Multiply the second identity in (HTW.13) by \(d\),
expand the alternating sum, and use the first identity:
only \(w=1\) survives and the result is \(df\).
Cancellation is valid in the integral universal ring,
so the identity is integral.  Applying \(w\) to the
identity for \(w^{-1}\) proves the interchanged version.
This construction inverts neither \(n!\) nor \(d\)
in the coefficient ring of any module.
\end{proof}
% END HTW10

% BEGIN HTW11
\paragraph{HTW11. The actual Iwahori trace identities.}
\label{claim:HTW11}
Let \(E_j:W\to W\) be the finite trace from Iwahori
fixed vectors to the maximal compact fixed vectors
at place \(j\), followed by inclusion.  It commutes
with symmetric polynomials at \(j\), and
\[
 \sum_T a_T(U_j)E_jb_T(U_j)=N_{Q_j}(U_j).
 \tag{HTW.14}
\]
If \(v\) is maximal-compact fixed at \(j\), then
\[
 E_j z(U_j)v=\partial(N_{Q_j}z)(U_j)v .
 \tag{HTW.15}
\]
Traces and operators at different places commute and
preserve the other maximal-compact fixed spaces.
\begin{proof}
Here are the finite Hecke calculations behind the
identities.  Bruhat decomposition of the reduction of
the maximal compact gives one Iwahori double coset
for each \(w\in S_n\); its trace is \(E=\sum_wT_w\).
For an adjacent transposition \(s=(i,i+1)\), enumerating
the root-subgroup representatives gives
\[
 \begin{split}
 T_s^2&=(Q-1)T_s+Q,\\
 T_sX_i&=X_{i+1}T_s+(Q-1)X_i,\\
 T_sX_{i+1}&=X_iT_s-(Q-1)X_i .
 \end{split}
 \tag{HTW.16}
\]
It commutes with the other \(X_k\).  In the first
enumeration one term returns to the identity and the
remaining nonzero residue classes give \(T_s\);
inserting the diagonal translation before or after
these representatives gives the last two formulas.
The scalar normalization \(U_i=Q^i\Theta(e_i(x))\)
is the one making them (HTW.16).  Reduced products
give \(T_w\); a length-increasing product enumerates
each corresponding Bruhat cell once.

Over the fraction field in the \(X_i\), these relations
are represented by
\[
 T_sf=Q\,s(f)+(Q-1)X_i
                        \frac{f-s(f)}{X_i-X_{i+1}}.
 \tag{HTW.17}
\]
The numerator is divisible by the displayed difference.
This representation is faithful on the Bernstein span:
the highest permutation term of \(T_w\) has nonzero
coefficient, and the distinct permutations are linearly
independent as automorphisms of the rational function
field.  The latter follows by the shortest-relation
argument: translate a supposed relation by multiplication
by a variable on which two automorphisms differ and
subtract.  The double-coset basis realizes this same
span in the actual Hecke algebra; its distinct
translation and Bruhat terms have disjoint leading
double-coset support.

Summing (HTW.17) over reduced Bruhat products gives
\[
 Ef=\sum_w\frac{w(N_Qf)}{w(d)}
                    =\partial(N_Qf).
 \tag{HTW.18}
\]
For completeness, its coefficient at the longest
permutation is the product of the reversed adjacent
factors \((X_i-QX_k)/(X_i-X_k)\).  The relations
\(T_sE=QE\) make the output symmetric, so permuting
this coefficient determines all the other coefficients.
This proves (HTW.18) and also verifies the normalization
at \(n=1\), where both products are \(1\).

Now HTW10 applied inside (HTW.18) gives
\(\sum_Ta_TEb_T=N_Q\) as operators on the faithful
model, hence in the Hecke algebra.  This proves
(HTW.14) on every actual module.  On a vector satisfying
\(T_sv=Qv\), the same calculation with (HTW.16)
gives \(Ezv=\partial(N_Qz)v\), proving (HTW.15).
All identities are polynomial identities over the
integral coefficient ring with \(Q\) inverted;
their verification in its fraction field therefore
does not invert any integer in the actual module.
Different places act in different factors of the
product group.  Finite sums and equivariance give the
remaining commutations and preservation assertions.
\end{proof}
% END HTW11

% BEGIN HTW12
\paragraph{HTW12. The two maps before completion.}
\label{claim:HTW12}
Let \(S\) be the simultaneous ordered splitting algebra
of the monic polynomials \(P_j\) over \(B\), with roots
\(u_{ji}\).  There are polynomial-semilinear maps
\[
 \phi:S\otimes_BM_0\longrightarrow W,\qquad
 \psi:W\longrightarrow S\otimes_BM_0
\]
whose two compositions are multiplication by
\(f_0=\prod_jN_{Q_j}(U_j)\), evaluated on the respective
sides.
\begin{proof}
Write \(\mathrm{inc}:M_0\hookrightarrow W\) for the
original inclusion through the level equivalence.
The product of all \(E_j\) lands in its image; define
\(\mathrm{tr}\) by
\(\mathrm{inc}\,\mathrm{tr}=\prod_jE_j\).
It is \(R\)-semilinear through \(\beta\).  HTW9 makes
the formula
\(\phi(z(u)\otimes m)=z(U)\mathrm{inc}(m)\)
descend through the splitting relations.

For a choice \(\mathbf T=(T_j)_j\) put
\(a_{\mathbf T}=\prod_ja_{T_j}(U_j)\) and similarly
for \(b_{\mathbf T}\).  Define the reverse map explicitly:
\[
 \psi(v)=\sum_{\mathbf T}a_{\mathbf T}(u)\otimes
                                     \mathrm{tr}(b_{\mathbf T}(U)v).
 \tag{HTW.19}
\]
The divided-difference pairing is linear over symmetric
polynomials.  HTW10's two dual expansions show
\(\sum a_T\otimes b_T\) is balanced with respect to
multiplication by any variable over the symmetric
coefficient ring: pair its two sides against arbitrary
polynomials by \(\partial\); the two dual expansions
recover equality.  Equivalently the expansions exhibit
this pairing as a perfect finite-projective pairing.
HTW11's formula, with HTW9 for symmetric scalars, gives
the same balancing for (HTW.19).  Thus \(\psi\) is
semilinear in every \(U_{ji}\), as well as in \(R\).

For \(\phi\psi\), substitute (HTW.19), move operators
at distinct places past one another, and use
(HTW.14) at each place.  The result is \(f_0(U)\).
For \(\psi\phi\), use (HTW.15) successively at every
place; (HTW.13) then gives \(f_0(u)\).
The finite product argument also covers \(q=0\),
when the splitting algebra is \(B\), the trace and
inclusion are inverse, and \(f_0=1\).
\end{proof}
% END HTW12

% BEGIN HTW13
\paragraph{HTW13. Enlargement to the literal correction polynomial.}
\label{claim:HTW13}
The maps in HTW12 can be chosen with both compositions
equal to \(F\) in (HTW.2).
\begin{proof}
Each factor \(X_k-Q_jX_i\), \(i<k\), in
\(N_{Q_j}\) is a factor of \(F\), namely the second
factor of its ordered pair \((k,i)\).
Thus \(F=f_0c\) for an integral polynomial \(c\).
Keep \(\phi\) and replace \(\psi\) by
\(\psi\circ c(U)\).  Polynomial semilinearity and
commutation give
\(\phi\psi=F(U)\), \(\psi\phi=F(u)\).
No genericity, localization, or cancellation of \(F\)
is used.  This is precisely the original
\(\mathrm{TwoMaps}\) statement.
\end{proof}
% END HTW13

% BEGIN HTW14
\paragraph{HTW14. The complete original operator-module record.}
\label{claim:HTW14}
An actual \(\mathrm{LevMod}\) with
\(\mathrm{LocalLevelConditions}\) and
\(Q_j\equiv1\bmod\ell^N\) yields the original
\(\mathrm{OpMod}\), together with both fields of
\(\mathrm{OpRel}\) and its two-map field, with exponent \(n!\).
\begin{proof}
Keep the actual coefficient ring \(R\), residual datum,
pair \(p\), coefficient map \(\beta\), and the original
weight-variable identities of \(\mathrm{LevMod}\).
Choose \(M_1,W\), the two polynomial actions, and trace
as in HTW2.  The finite-level module field makes \(M_1\)
finite over \(\Lambda\).  The projectivity field,
transported through the concrete quotient (HTW.4),
gives the required \(\Lambda[\Delta]\)-projectivity.
Formula (HTW.5) gives the exact generator relation
\(\gamma_{ji}=1+Y_{ji}\).  HTW1 supplies \(a_j\ge N\),
trace surjectivity, and its exact kernel.
HTW2 supplies the \(U\)-intertwining.
These are all fields of \(\mathrm{OpMod}\).
HTW13 gives \(\mathrm{OpRel.two}\); HTW8 gives its
two universally quantified trace relations, including
every \(g\) in the original group.  Thus no operator
field remains an additional assumption.
\end{proof}
% END HTW14

\subsection{The completed coefficient ring and free diamond module}

% BEGIN HTW15
\paragraph{HTW15. A canonical simultaneous topologically nilpotent summand.}
\label{claim:HTW15}
On a finite module over complete noetherian local \(\Lambda\),
a finite commuting family \(T_s\) has a direct summand
on which all \(T_s\) are topologically nilpotent.
Its projector commutes with every endomorphism commuting
with the family.
\begin{proof}
Let \(C\) be the commutative \(\Lambda\)-subalgebra generated
by the operators in the finite module
\(\operatorname{End}_\Lambda(M)\).  It is finite over
\(\Lambda\): each generator is integral by the determinant
trick on a finite module, and there are finitely many
generators.  It is complete.  The Artinian algebra
\(C/\mathfrak m_\Lambda C\) is a product of local
Artinian algebras.  Sum the primitive factor idempotents
in which every \(T_s\) is in the maximal ideal.
Idempotents lift uniquely through the complete
\(\mathfrak m_\Lambda\)-adic ring: lift successively
through nilpotent quotients, where the equation
\(e^2-e=0\) has invertible derivative \(2e-1\), and
take the inverse limit.

On the selected summand sufficiently high powers of
each \(T_s\) lie in \(\mathfrak m_\Lambda C\), hence
their powers tend to zero.  On any omitted factor
one \(T_s\) is a unit; if its powers on a vector tend
to zero, multiplication by its inverse and separatedness
of the finite module force that vector to be zero.
The selected range is therefore exactly the
simultaneous topologically nilpotent submodule.
The projector belongs to \(C\), so any endomorphism
centralizing the \(T_s\) centralizes it.  Inclusion
and projection are split \(\Lambda\)-linear maps.
\end{proof}
% END HTW15

% BEGIN HTW16
\paragraph{HTW16. Power-series actions and their semilinear maps.}
\label{claim:HTW16}
On the summand in HTW15 the commuting operators extend
to an action of \(R[[X_s]]\), with \(X_s\) acting as \(T_s\),
when the original operators and the projector commute
with \(R\).  A map between two such finite
\(\Lambda\)-modules that intertwines coefficients and
variables intertwines every power series.
\begin{proof}
Modulo \(\mathfrak m_\Lambda^k\), some power of every
\(T_s\) is zero, uniformly on the finite module.
Since there are finitely many variables, all monomials
of sufficiently large total degree act as zero.
Evaluate a power series by its finite truncation in
this quotient.  These evaluations are compatible as
\(k\) varies, so completeness gives a unique action.
Addition and multiplication can be checked in each
finite truncation, where they are polynomial identities.

A coefficient-semilinear additive map is
\(\Lambda\)-linear when the two coefficient maps agree
on \(\Lambda\), and hence preserves
\(\mathfrak m_\Lambda^k\)-multiples.  In every quotient
its variable-intertwining property extends to the
finite polynomial truncation.  Passing to the inverse
limit proves the full power-series assertion.
Thus no unproved continuity of an arbitrary additive
map is being used.
\end{proof}
% END HTW16

% BEGIN HTW17
\paragraph{HTW17. Freeness over the exact diamond quotient.}
\label{claim:HTW17}
Apply HTW15 to \(M_1\) with the operators
\(U_{ji}-u^0_{ji}\) and \(Y_{ji}\), where \(u^0_{ji}\in O\)
are any chosen lifts.  Let its range be \(M\).
Then for some finite \(r\),
\[
 M\simeq\left(
  \Lambda[[Y_{ji}]]/
       ((1+Y_{ji})^{\ell^{a_j}}-1)_{j,i}
                   \right)^r
 \tag{HTW.20}
\]
as a module for the displayed power-series action.
\begin{proof}
The projector commutes with \(\Delta\), since
\(\gamma_{ji}=1+Y_{ji}\) and all polynomial operators
commute.  Therefore \(M\) is a direct summand of
the projective \(\Lambda[\Delta]\)-module \(M_1\).
It is finite.  The group algebra \(\Lambda[\Delta]\)
is local: modulo \(\mathfrak m_\Lambda\) its augmentation
ideal is nilpotent, because the residue characteristic
is \(\ell\) and each generator has \(\ell\)-power order.
Its only simple quotient is the residue field.
A finite projective module over a local ring is free:
lift a residue-field basis, obtain a surjection from
a finite free module by Nakayama, split it by
projectivity, and use Nakayama again to kill the
complement.

Finally \(Y_{ji}\mapsto\gamma_{ji}-1\) identifies
the ring in (HTW.20) with \(\Lambda[\Delta]\).
Polynomial division by the monic relations gives
the basis \(\prod Y_{ji}^{b_{ji}}\), \(0\le b_{ji}<\ell^{a_j}\).
The relations are distinguished, so the same division
holds for power series over the complete local ring,
by successive reduction modulo powers of its maximal
ideal.  Both sides are finite complete free
\(\Lambda\)-modules with this basis, and the generator
map identifies them.  It respects the action on \(M\)
by HTW16 and the exact generator relation.
\end{proof}
% END HTW17

% BEGIN HTW18
\paragraph{HTW18. The augmentation quotient is the correct primary summand.}
\label{claim:HTW18}
The trace identifies
\[
 X=M/(Y_{ji})M
       \simeq W_{u^0},
 \tag{HTW.21}
\]
where the right side is the simultaneous topologically
nilpotent summand for \(U_{ji}-u^0_{ji}\) on \(W\).
The isomorphism and its split inclusion into \(W\)
intertwine all polynomial \(U\)-operators.
\begin{proof}
HTW1 gives \(W=M_1/(Y_{ji})M_1\).  The projector
in HTW15 commutes with every \(Y_{ji}\), so taking
this quotient splits as the quotient of its range
and the quotient of its complementary range.
On \(W\) all \(Y_{ji}\) are zero.  The image of the
projector is exactly the summand where the remaining
\(U_{ji}-u^0_{ji}\) are topologically nilpotent:
this follows either from its local-factor construction
or from uniqueness of the decomposition in HTW15.
Thus the trace induces the stated isomorphism and
the splitting.  Commutation with the \(U\)-operators
gives all the intertwining statements.  In particular
one has both an inclusion \(j:X\to W\) and a
projection \(\kappa:W\to X\), with \(\kappa j=1\).
\end{proof}
% END HTW18

% BEGIN HTW19
\paragraph{HTW19. Completion of the ordered splitting algebra.}
\label{claim:HTW19}
The source ring \(\mathrm{Sp}(P,\beta u^0)\) is the
factor of the finite ordered splitting algebra \(S\)
on which every \(u_{ji}-\beta(u^0_{ji})\) is
topologically nilpotent.  The maps of HTW13 restrict
to
\[
 \mathrm{Sp}(P,\beta u^0)\otimes_BM_0
          \ \rightleftarrows\ W_{u^0},
 \tag{HTW.22}
\]
with both compositions still equal to the literal \(F\).
\begin{proof}
A monic polynomial can be split successively by
adjoining one root, dividing by its monic linear
factor, and repeating.  Thus \(S\) is finite over
complete local \(B\).  Its quotient modulo
\(\mathfrak m_B\) is Artinian.  Select the factors
where all the ordered root differences have
residue zero, and lift their summed idempotent as
in HTW15.  If no such factor exists the selected
ring is the zero ring, which is allowed in this
description.

The defining power-series quotient for
\(\mathrm{Sp}\) is precisely this factor.  In one
direction, its topologically nilpotent root
differences give an evaluation of the power series
and satisfy every coefficient of the splitting
relations.  In the other, those relations give an
ordered splitting of each \(P_j\).  The universal
maps are inverse on coefficients and root
differences and hence, by completeness, on the
whole rings.  The factor is a finite direct summand
of \(S\).

Tensoring its idempotent with \(M_0\) preserves this
direct-summand description, without any flatness
assumption on \(M_0\).  The two maps of HTW13
intertwine every root difference, so they carry
the chosen summands into one another.  Restriction
of their already proved composition identities
gives (HTW.22).  These modules are finite over
\(\Lambda\): the left side is a quotient of a
finite direct sum of copies of \(M_0\), and the
right side is a summand of the finite quotient \(W\).
Their root-difference actions are topologically
nilpotent also in the \(\Lambda\)-adic topology.
If \(M_0=0\), the left module is zero and this
assertion is immediate there; the right summand
already has its \(\Lambda\)-adic nilpotence by
definition.  Otherwise the image of \(B\) in the endomorphisms
of \(M_0\) is a finite local \(\Lambda\)-algebra:
it is a \(\Lambda\)-submodule of a finite
endomorphism module and a quotient of local \(B\).
Its maximal ideal is the radical of
\(\mathfrak m_\Lambda\) times that image.
Consequently powers of \(\mathfrak m_B\)
act through arbitrarily high powers of
\(\mathfrak m_\Lambda\).  The same argument
on the finite splitting-algebra tensor module
converts nilpotence of each root difference
modulo \(\mathfrak m_B\) into the required
\(\Lambda\)-adic topological nilpotence.
\end{proof}
% END HTW19

% BEGIN HTW20
\paragraph{HTW20. The completed ring and the exact coefficient-map identity.}
\label{claim:HTW20}
From the \(\mathrm{OpMod}\) of HTW14 define
\[
 R'=R[[X_{ji},Y_{ji}]],\qquad U'_{ji}=u^0_{ji}+X_{ji}.
 \tag{HTW.23}
\]
Extend its pair by the constant coefficient inclusion.
There are original maps
\[
 s:\Lambda[[Y_{ji}]]\longrightarrow R',
 \quad
 \theta:R'\longrightarrow\mathrm{Sp}(P,\beta u^0)
\]
with the original weight and pair identities, and
\[
 \theta(U'_{ji})=u_{ji},\qquad
 \theta(s(z))=
       B\longrightarrow\mathrm{Sp}\bigl(\lambda(z|_{Y=0})\bigr).
 \tag{HTW.24}
\]
\begin{proof}
First \(\beta\) agrees with \(\lambda\) on
\(\Lambda\).  On a weight variable both values are
the prescribed scalar times the corresponding
ordinary character minus \(1\), by the
\(\mathrm{OpMod}\) weight and pair fields.
They agree on \(O\).  These are local maps of
complete coefficient rings, hence adically
continuous, so their agreement on polynomials
gives agreement on power series.

Define \(s\) by the original weight map on
coefficients and by the right-hand \(Y\)-variables.
Define \(\theta\) on \(R\) through \(\beta\),
on \(X_{ji}\) by \(u_{ji}-\beta(u^0_{ji})\),
and on every \(Y_{ji}\) by zero.  These variable
images are topologically nilpotent, so the
complete power-series maps exist.  Formula
(HTW.24) holds on coefficients and variables,
hence on all power series.

Constant base change preserves every field of
the full pair, and \(\beta.\mathrm{pair}\) gives
the full determinant and ordinary-character
identity after \(\theta\).  This is equality of
polynomial laws over every scalar algebra.
The prescribed weight equality for \(s\) follows
from that for \(R\), and
\(U'_{ji}-u^0_{ji}=X_{ji}\in\mathfrak m_{R'}\).
The ring \(R'\) is still an original complete
noetherian local coefficient ring with the same
residue field.
\end{proof}
% END HTW20

% BEGIN HTW21
\paragraph{HTW21. All fields of TWRing and TWMaps.}
\label{claim:HTW21}
HTW14 produces the literal \(\mathrm{TWRing}\)
and \(\mathrm{TWMaps}\) with exponent \(n!\), for
any specified \(u^0\), with the same \(b,J,N\).
\begin{proof}
Use the ring, pair, map \(s\), and roots in HTW20.
Let \(R'\) act on the module \(M\) in HTW17:
the left \(X\)-variable acts by \(U-u^0\) and
the right \(Y\)-variable by the original \(Y\).
HTW16 constructs this action.  The isomorphism
(HTW.20), with the same \(a_j\ge N\), is exactly
the required free diamond-quotient model for
\(\mathrm{TWRing}\).

Compose the two maps of (HTW.22) with the inverse
isomorphisms in (HTW.21).  Polynomial
semilinearity holds by construction.  HTW16
extends it to every \(t\in R'\), so these are
the exact \(\phi,\psi\) fields of \(\mathrm{TWMaps}\).
Their compositions are \(F(U')\) and its
\(\theta\)-image, without removing any factor.
HTW20 supplies \(\theta.\mathrm{pair}\),
\(\theta(U')\), and the coefficient-map identity.

Finally polynomial evaluation
\(U\mapsto u^0+X,\ Y\mapsto Y\) carries the
trace polynomials in (HTW.10) to the corresponding
trace polynomials of the constantly extended
pair.  This follows from full polynomial-law
naturality, including its trace coefficient.
Inclusion \(M\hookrightarrow M_1\) is injective
and intertwines this evaluation.  Thus both
\(n!\)-th power relations hold on every \(m\in M\),
for every \(j,i,g\).  This checks the last two
fields, completing every field of both records.
\end{proof}
% END HTW21

% BEGIN HTW22
\paragraph{HTW22. Transport through the original group isomorphism.}
\label{claim:HTW22}
The \(\mathrm{OpMod}\) and all its relations transport
through the original injective map \(e:G\to G'\)
of global data, when its map to the fixed base
commutes.  The coefficient rings, modules, roots,
base \(b\), ideal \(J\), and exponent do not change.
\begin{proof}
Such a global-data map is already surjective by
its defining fields.  With the stated injectivity,
it is an isomorphism; a continuous bijection
between these compact Hausdorff groups has a
continuous inverse.  Use that inverse \(k\) to
pull the residual datum and full determinant law
to \(G'\).  On the local groups the structural
maps intertwine the prescribed characters,
polarization, ordinary relations and type kernels;
thus the pulled object is the original pair.
For the type kernels specifically, a point lies
in the transported inertia image exactly when
its inverse image lies in the old one, and each
residual-character kernel is preserved by the
same character equality.  This checks their
intersection, not just the ambient inertia group.

The base compatibility says \(h\circ k=h'\);
hence the old coefficient map still sends the
pair to the fixed base pair.  All weight
characters remain the same.  In the trace
relations replace \(g'\) by \(k(g')\);
the group-algebra map sends \([e(f_j)]\) to
\([f_j]\) and similarly for \(t_j\).
Polynomial-law pullback commutes with coefficient
base change and trace, so both relations are
exactly the old ones.  All module and two-map
fields are unchanged.  This proves the literal
\(\mathrm{opTransport}\) consumer.
\end{proof}
% END HTW22

\subsection{One base and one exponent for the entire auxiliary family}

% BEGIN HTW23
\paragraph{HTW23. The full original HasAutOfAlt implication.}
\label{claim:HTW23}
Under the original \(\mathrm{LocPAlt}\), if
\(V_0.S'=s.S\) and \(\mathrm{LevelAlt}(B,J)\),
then \(\mathrm{HasAut}(B,J)\).  The same one
base \(b\), one residual-root lift \(u^0\), and
one exponent \(n!\) work for every auxiliary
integer \(N\).
\begin{proof}
Choose \(b\) once from \(\mathrm{LevelAlt}\),
before any auxiliary level.  Now fix an arbitrary
original Taylor--Wiles group family \(T\) and its
place realization \(Q\).  For each \((j,i)\)
choose \(u^0_{ji}\in O\) lifting the residual
character value \(\overline\chi_i(T.sg_j)\);
the residue map is surjective.  This choice
precedes \(N\).

For each \(N\), use the given \(Q.V_N\),
places \(Q.v_N,Q.w_N\), and prolongations \(Q.B_N\).
Their distinctness, avoidance, exact set of
auxiliary places, splitting, and compatibility
with the base map are fields of the same
original realization.  Thus the single
\(\mathrm{LevelAlt}\) statement supplies the
actual \(\mathrm{LevMod}\) and its separating
classical family at that level.  Avoidance of
\(s.S\), which contains the places above \(\ell\),
shows the auxiliary \(w_j\) are away from
\(\ell\), by restriction of their prolongations.

Apply HTW6 to the given \(Q.fr_N,Q.\tau_N\).
Its cardinality equality identifies the actual
\(Q_j\) with \(T.Nq_N\), and its comparison uses
the full \(\mathrm{LevPtAlt.lg}\) of the same
classical points.  The congruence field of \(T\)
then gives the hypothesis of HTW14.
Transport its operator module along the actual
\(Q.e_N\) using HTW22.  The Frobenius and tame
compatibility fields of \(Q\) identify the
transported elements with \(T.ph_N,T.ta_N\).
Apply HTW21 with the already chosen \(u^0\).

Choose these resulting rings and maps for every
\(N\).  Together with \(b,u^0\), their residual
equalities and the fixed exponent \(n!\), they
are exactly \(\mathrm{TWAut}(T)\).  Since \(T,Q\)
were arbitrary, this is the original
\(\mathrm{HasAut}\).  Empty auxiliary families
are included by the empty products already
checked above.  The abstract implication does
not impose a new lower bound on \(n\); its
zero-rank matrix and trace assertions were
included explicitly.
\end{proof}
% END HTW23

% BEGIN HTW24
\paragraph{HTW24. Actual same-ideal automorphy from the Hida level construction.}
\label{claim:HTW24}
For every original \(n\ge1\) datum covered by HLM26,
its actual \(\mathrm{IsJ}(B,K,J)\), and the original
local published hypotheses, the full
\(\mathrm{HasAut}(B,J)\) holds with exactly that ideal.
\begin{proof}
HLM26 proves the literal all-level
\(\mathrm{LevelAlt}\) producer from the supplied
unitary Hecke datum and \(\mathrm{IsJ}\).  In
particular it constructs one \(b\) satisfying
(HTW.1) before all \(q,V,h,v,w,B_q\), rather
than producing an unrelated base at each level.
Apply HTW23 to this output and to
\(V_0.S'=s.S\).  Every choice in that proof
keeps \(b\) and hence the same support ideal.
The resulting \(\mathrm{HasAut}\) is the full
original record with both trace relations,
the exact coefficient map and correction
polynomial, for every original auxiliary
family.  Nothing excludes \(\ell=3\).

This completes the level-to-automorphy producer.
It does not itself install both rank-four
\(\mathrm{Link}\) objects or identify the
entire selected \(\mathrm{COUNT\_II\_SUPPORT}\)
contract with this one endpoint.  Those wider
contracts retain their other original fields.
\end{proof}
% END HTW24

\subsection{Scope of this completion}

The three published fields of \(\mathrm{LocPAlt}\)
remain explicitly retained at their original scope.
The proof above supplies the intervening finite
Hecke, matrix, determinant-law and complete-ring
arguments and constructs the full output record.
It preserves the fixed coefficient DVR, the
given complex comparison, the auxiliary-place
realization, the complete ordinary pair, the
same \(J\), the exact ordered correction factor,
and all integer \(N\).  In particular the
nilpotence exponent is uniform because it
comes from one local Jacquet dimension bound,
while the existence of the free module comes
from projectivity over the finite diamond
group algebra, not from dividing by its order.
